% AUTO-GENERATED by codes-for-paper/outflow_composition/tables/stage2_to_latex.py
% Do not edit by hand — regenerate from the CSV instead.
\begin{table*}[!ht]
\centering
\caption{Minimum magnetisation following \citet{ZhangPeer2009}: the smallest $1+\sigma$ for which the photospheric blackbody of the fireball stays below the best-fit spectrum across the test band, for the episodes whose best-fit model has no blackbody.
The test band is the observer-frame \SI{10}{\kilo\electronvolt}--\SI{10}{\mega\electronvolt}, the gamma-ray luminosity $L_\gamma$ is integrated over the observer-frame \SI{10}{\kilo\electronvolt}--\SI{10}{\giga\electronvolt}, and $L_\mathrm{w}=L_\gamma$.
The base radius is $R_0=c\,t_\mathrm{v}/(1+z)$ and the branch is $\eta\geq\eta_*$.
Values are quoted at the spectroscopic redshift $z=4.35$ for \grbzeroeightzeroninesixteenC\ and at the assumed redshifts $z=1,2,3,5,7$ for the other bursts.
Errors are asymmetric $1\sigma$ intervals from $10000$ Monte Carlo samples (seed $1962439590$).}
\label{tab:outflow_stage2}
\resizebox{\textwidth}{!}{
\begin{threeparttable}
\renewcommand{\arraystretch}{1.25}
\begin{tabular}{llcccccc}
\toprule
Episode & Model & $z=4.35$ & $z=1$ & $z=2$ & $z=3$ & $z=5$ & $z=7$ \\
\midrule
\multicolumn{8}{l}{\textbf{\grbzeroeightzeroninesixteenC}} \\
TR2 & \band & $19.4_{-3.2}^{+2.4}$ &  &  &  &  &  \\
TR3 & \sbpl & $63.4_{-2.6}^{+2.8}$$^{\ddagger}$ &  &  &  &  &  \\
TR4$^{\dagger}$ & \sbpl & $44.0_{-10.0}^{+13.0}$ &  &  &  &  &  \\
EX1$^{\dagger}$ & \sbpl & $38.7_{-6.8}^{+7.5}$$^{\ddagger}$ &  &  &  &  &  \\
TR5$^{\dagger}$ & \sbpl & $24.8_{-5.3}^{+6.7}$ &  &  &  &  &  \\
\midrule
\multicolumn{8}{l}{\textbf{\grbfourteenzerotwozerosixB}} \\
T90$^{\dagger}$ & \band &  & $1.727_{-0.039}^{+0.040}$$^{\ddagger}$ & $2.940_{-0.066}^{+0.069}$$^{\ddagger}$ & $3.700_{-0.083}^{+0.086}$$^{\ddagger}$ & $4.60_{-0.10}^{+0.11}$$^{\ddagger}$ & $5.13_{-0.12}^{+0.12}$$^{\ddagger}$ \\
TR2 & \band &  & $25.197_{-1.513}^{+1.833}$ & $17.752_{-1.034}^{+1.265}$ & $15.201_{-0.865}^{+1.075}$ & $13.09_{-0.74}^{+0.92}$ & $12.13_{-0.68}^{+0.85}$ \\
TR3 & \band &  & $6.863_{-0.948}^{+1.092}$ & $5.920_{-0.777}^{+0.905}$ & $5.562_{-0.708}^{+0.836}$ & $5.25_{-0.64}^{+0.78}$ & $5.10_{-0.61}^{+0.75}$ \\
TR4 & \band &  & $8.624_{-0.452}^{+0.563}$ & $7.295_{-0.372}^{+0.467}$ & $6.776_{-0.340}^{+0.431}$ & $6.31_{-0.31}^{+0.40}$ & $6.09_{-0.30}^{+0.38}$ \\
TR5$^{\dagger}$ & \pl &  & $1.087_{-0.087}^{+0.201}$$^{\ddagger}$ & $1.679_{-0.235}^{+0.314}$$^{\ddagger}$ & $2.005_{-0.283}^{+0.378}$$^{\ddagger}$ & $2.35_{-0.33}^{+0.45}$$^{\ddagger}$ & $2.54_{-0.36}^{+0.48}$$^{\ddagger}$ \\
TR6 & \cpl &  & $7.843_{-0.566}^{+0.602}$$^{\ddagger}$ & $10.845_{-0.785}^{+0.829}$$^{\ddagger}$ & $12.190_{-0.885}^{+0.932}$$^{\ddagger}$ & $13.38_{-0.97}^{+1.03}$$^{\ddagger}$ & $13.89_{-1.02}^{+1.07}$$^{\ddagger}$ \\
\midrule
\multicolumn{8}{l}{\textbf{\grbtwentythreeeleventwentynineC}} \\
EX1 & \sbpl &  & $5.235_{-0.607}^{+0.477}$ & $3.508_{-0.366}^{+0.290}$ & $2.975_{-0.279}^{+0.226}$ & $2.58_{-0.20}^{+0.17}$ & $2.42_{-0.16}^{+0.14}$ \\
TR2 & \sbpl &  & $4.791_{-0.636}^{+0.485}$ & $3.081_{-0.348}^{+0.277}$ & $2.591_{-0.244}^{+0.207}$ & $2.26_{-0.15}^{+0.14}$ & $2.14_{-0.11}^{+0.11}$ \\
\bottomrule
\end{tabular}
\begin{tablenotes}
\footnotesize
\item[$\dagger$] No pulse-based variability timescale is available and the episode duration stands in for $t_\mathrm{v}$. Unlike $\Gamma_{\min}$, $1+\sigma_\mathrm{min}$ is not monotonic in $t_\mathrm{v}$, so the entry may lie above or below its value for the true timescale.
\item[$\ddagger$] The binding blackbody-to-spectrum ratio lies at the lower edge of the test band (\SI{10}{\kilo\electronvolt}) in at least \SI{50}{\percent} of the Monte Carlo draws.
\end{tablenotes}
\end{threeparttable}
}
\end{table*}
